\subsection{贫集}

定义 2. 拓扑空间 X 的子集 A 称为贫集，如果它是一个可数无处稠密集族的并。

等价地，A 是贫集，如果它含于可数多个闭集的并，而这些闭集中的每一个都没有内点。

贫集完全可能在 X 中稠密；甚至整个空间 X 本身就可以是贫集。

后一种情形的例子由任何没有孤立点的可数豪斯多夫空间给出，例如有理数直线 Q 。但作为 X 中的贫集的拓扑空间 X 未必是可数的（见习题 9）。

空间 X 中贫集的每个子集都是贫集，且可数贫集族之并是贫集。

设 Y 是 X 的子空间。Y 的子集 A 称为相对于 Y 的贫集，如果 A 作为拓扑空间 Y 的子集来看是贫集。由第 1 号命题 2 推出：若 A 是 Y 的子集且相对于 Y 是贫集，则 A 相对于 X 是贫集；又若 Y 在 X 中开，则 Y 的每个相对于 X 为贫集的子集 A 相对于 Y 也是贫集。

